%% Cover Letter
%% Marketing Letters — Replication Corner
%% Manuscript: "Does AI Reduce the Reputation Premium in Knowledge Diffusion?
%%              A Conceptual Replication Using Chinese Journal Articles"

\documentclass[12pt,a4paper]{article}

\usepackage[top=2.54cm,bottom=2.54cm,left=3.0cm,right=3.0cm]{geometry}
\usepackage[utf8]{inputenc}
\usepackage[T1]{fontenc}
\usepackage{times}
\usepackage{setspace}
\usepackage{hyperref}
\usepackage{microtype}
\usepackage{parskip}

\singlespacing
\setlength{\parskip}{0.9em}

\begin{document}

\noindent
[Author names and affiliation --- blinded for review]

\medskip

\noindent
\today

\bigskip

\noindent
\textit{Marketing Letters} Editorial Office\\
Editor-in-Chief

\bigskip

\noindent
Dear Editor,

\bigskip

We submit for your consideration a manuscript entitled \textbf{``Does AI
Reduce the Reputation Premium in Knowledge Diffusion? A Conceptual
Replication Using Chinese Journal Articles''} for the \textit{Replication
Corner} section of \textit{Marketing Letters}.

\medskip
\noindent\textbf{Motivation and fit.}
Jo and Li (2026, \textit{Marketing Letters}) provide the first direct
evidence that AI-assisted writing attenuates the institutional reputation
premium in academic knowledge diffusion, using a sample of SSRN marketing
working papers.
Their finding carries broad implications for the sociology of knowledge and
academic publishing, but it rests on a single context: English-language
preprints with no editorial gatekeeping.
A key open question is whether the effect generalizes to peer-reviewed
publication venues, and to non-Western academic contexts with distinct
institutional hierarchies.
Our study directly addresses this question and constitutes a natural
contribution to \textit{Marketing Letters}' own emerging conversation on AI
and knowledge diffusion.

\medskip
\noindent\textbf{What we do.}
We conceptually replicate Jo and Li (2026) using 480 peer-reviewed Chinese
business and management journal articles from CNKI (China National Knowledge
Infrastructure; 2023--2026).
The AI involvement measure is an abstract-level stylometric score generated
via a structured ChatGPT evaluation against a pre-specified dictionary of
23 AI writing-style markers, developed in turn from the taxonomy of Liang
et al.\ (2024, \textit{NEJM AI}) and validated on a 30-article
human-coded subsample (Pearson $r = 0.82$).
Institutional prestige is operationalized via a four-tier classification
based on China's C9, 985, and 211/Double First-Class designations,
coded by two trained research assistants ($\kappa = 0.93$).
We estimate OLS models with HC3-robust standard errors for two diffusion
outcomes (log downloads, log citations) across four specifications (binary
and continuous prestige operationalizations).

\medskip
\noindent\textbf{Key findings.}
Both main effects replicate robustly: institutional prestige
($\hat\beta = 0.522$, $p < .001$) and AI-style writing
($\hat\beta = 0.266$, $p = .002$) each independently predict greater
diffusion.
However, the critical Prestige\,$\times$\,AI-Style interaction is
consistently near zero and statistically insignificant across all four
specifications (point estimates ranging from $-0.085$ to $+0.013$;
all $p > .28$; 95\,\% CIs spanning zero in both directions).
This finding suggests that the reputation-equalizing potential of AI-style
writing documented for SSRN preprints may not extend to peer-reviewed
publication contexts where editorial gatekeeping provides an independent
quality signal---at least within Chinese business and management
scholarship.

\medskip
\noindent\textbf{Distinctive contributions.}
First, this is, to our knowledge, the first conceptual replication of Jo
and Li (2026), appearing in the same journal that published the original
study.
Second, our Chinese peer-reviewed context introduces two theoretical levers
absent from SSRN: editorial certification as an independent quality signal,
and an institutionalized prestige hierarchy embedded in national funding
and promotion criteria.
Third, the positive replication of both main effects strengthens the
cross-context generalizability of those effects, while the null interaction
identifies a meaningful scope condition.
Fourth, the ChatGPT-assisted scoring methodology offers a scalable
approach to measuring AI-style writing in non-English corpora.

\medskip
\noindent\textbf{Practical considerations.}
The manuscript is approximately 3,400 words of running text (plus
$\approx 500$ words of references), within the 4,000-word limit for
Replication Corner submissions.
It includes exactly three manuscript tables and a Web Appendix containing
full regression outputs, data construction details (including AI-style
score prompt, marker dictionary, and prestige coding rubric), robustness
checks, and a side-by-side coefficient comparison with Jo and Li (2026).
All data were retrieved from publicly available CNKI records; analysis
code and the scoring prompt will be deposited in an open repository
upon acceptance.
This manuscript has not been submitted elsewhere and represents
entirely original work.

\bigskip

We thank the editors for considering this submission and look forward to
the review process.

\bigskip

\noindent Sincerely,

\bigskip\bigskip

\noindent [Author names --- blinded for review]

\end{document}
